\section{Bipartite witness graphs and nonnegative tournament gaps}
\label{sec:bipartite-witness-completion}

Let $\mathcal W$ be the graph whose vertices are the directional failure
sources of bad missing pairs. Join $x,y$ when they are opposite failure
witnesses for some bad pair. Its edges are original missing pairs.
The two directional failure sets of each bad pair span a complete
bipartite subgraph of $\mathcal W$.

\begin{lemma}[Completion along a bipartition]
\label{lem:bipartite-selected-source-completion}
Suppose $\mathcal W$ is bipartite with parts $L,R$. Orient each bad
missing pair in the direction whose failure sources belong to $L$,
and all good pairs canonically relative to protection, obtaining $T$.
For every vertex $z$, put $K_z=N_T^+(z)\setminus N_D^+(z)$. Then
\[
 N_T^{++}(z)\subseteq
 (N_D^{++}(z)\setminus K_z)\cup(\mathcal F(z)\setminus K_z),
 \qquad g_T(z)\ge g_D(z)+2|K_z|-t(z).
\]
If $z\notin L$ and $K_z=\varnothing$, its two neighbourhoods are
unchanged. If $t(z)=0$, the second-neighbourhood inclusion is the
exact equality $N_T^{++}(z)=N_D^{++}(z)\setminus K_z$.
\end{lemma}
\begin{proof}
The direction choice is well-defined: the complete bipartite graph
between the opposite failure sets makes each of these sets monochromatic,
in opposite colours. Every chosen bad direction therefore fails only
at vertices of $L$.

No added second arc through an original first neighbour reaches a
protected far head: a good direction is convenient, and every failure
head of a bad direction is unprotected. Through an added good first
arc, the canonical protection compatibility rules out such a head.
For an added bad first arc $z\to a$, choose an opposite witness
$r\in R$, so $r\to a$ and $z\in U_D(r)$.
If $y$ is protected from $z$, then $y\in U_D(r)$ as well.
An original $a\to y$ would force $a\to z$, contrary to the missing
pair $za$. If $ay$ is missing, its direction $a\to y$ fails at $r$.
A good pair must therefore be oriented $y\to a$; for a bad pair,
our choice avoiding $R$ also forces $y\to a$.
This checks all new first intermediates and all possible second arcs.

Thus every genuinely new far second head belongs to $\mathcal F(z)$.
Every added first neighbour was missing from $z$, hence lies in exactly
one of the disjoint sets $N_D^{++}(z)$ and $\mathcal F(z)$.
Removing the entire $K_z$ proves the count.
For $z\notin L$ with no new first neighbours, no chosen bad direction
can fail at $z$, so there are no new second heads. Original paths remain.
The exact assertion for $t(z)=0$ follows in the same way.
\end{proof}

\begin{lemma}[Only one selected failure source]
\label{lem:single-selected-source-completion}
If every bad missing pair can be oriented so that its chosen directional
failure sources lie in $\{c\}$, the preceding conclusion holds for
canonical good directions and these bad directions. In particular,
every $z\ne c$ with $t(z)\le1$ and $g_D(z)>0$ remains deficient in
the same completion, regardless of whether $K_z$ is empty.
\end{lemma}
\begin{proof}
All cross-witness edges are incident to $c$: the chosen failure set
of every bad pair is precisely $\{c\}$, and the opposite set excludes
$c$. Thus use $L=\{c\}$ in the previous lemma. If $K_z$ is empty its
gap is unchanged; otherwise $g_T(z)\ge g_D(z)+2-t(z)>0$.
\end{proof}

\begin{lemma}[The zero-gap set has no sink]
\label{lem:nonnegative-tournament-three-zeros}
Let $T$ be a finite nonempty tournament with no sink, and suppose
$g_T(z)\ge0$ for every vertex. Then its zero-gap vertices induce a
tournament with no sink. In particular there are at least three of them.
\end{lemma}
\begin{proof}
The unweighted tournament SNP theorem ensures a zero-gap vertex exists.
Suppose the zero-gap set $S$ has a sink $c$ in $T[S]$.
The absence of a sink in $T$ gives an outgoing neighbour $h$ of $c$
outside $S$. Assign weights $w(c)=5/4$, $w(h)=9/8$, and weight one
to all other vertices. The weighted gap at $c$ is $1/8$.
Every other zero-gap vertex points to $c$, so its weighted gap is at
least $1/4-1/8=1/8$. Every originally positive-gap vertex has integer
gap at least one, and its weighted gap is at least $1-1/4-1/8>0$.
All weighted gaps are therefore strictly positive, a contradiction to
the weighted tournament SNP theorem \cite[Proposition~2.1]{FidlerYuster2007}.

For this rational-weight use, that theorem follows directly from the
unweighted tournament theorem without assumptions about a weighted
median optimization: multiply the weights by eight, replace each vertex
by a transitive tournament of its integer weight, and orient between
blocks as in $T$. Every vertex in a block has external gap equal to
eight times the rational weighted original gap, and nonnegative
additional internal gap.
If all weighted gaps were positive, the expanded tournament would be
all-deficient, contradicting the unweighted theorem.
A tournament on at most two vertices always has a sink, proving the
last assertion.
\end{proof}
